% ============================================================
% Formulario: tabella riassuntiva delle distribuzioni notevoli (fine della Parte I)
% ============================================================
\clearpage
\argomento[T]{Tabella delle distribuzioni}{Le distribuzioni notevoli in un colpo d'occhio: il numero accanto al nome
  porta alla voce}[Tabella delle distribuzioni]
\colorlet{fcolore}{mspiform!80!black}\colorlet{fsfondo}{mspiform!5}

\newcommand{\nomedist}[3]{\textbf{#1}\ \vv{#2}\newline{\footnotesize\color{mspigray}#3}}

\vspace{2mm}
{\small
\renewcommand{\arraystretch}{2.1}
\noindent\begin{tabular}{P{38mm}P{22mm}P{54mm}P{18mm}P{21mm}}
  \toprule
  \textbf{Distribuzione} & \textbf{Valori} & \textbf{pmf o densità} & \textbf{Media} & \textbf{Varianza} \\
  \midrule
  \multicolumn{5}{l}{\color{fcolore}\Lato\bfseries Discrete} \\
  \rowcolor{fsfondo}
  \nomedist{Bernoulli}{bernoulli}{una prova: successo o insuccesso} & $\{0, 1\}$ & $p_X(1) = p$, \ $p_X(0) = 1 - p$
    & $p$ & $p\,(1 - p)$ \\
  \nomedist{Binomiale $\mathcal{B}(N, p)$}{distbinomiale}{successi in $N$ prove} & $0, 1, \ldots, N$ &
    $\dbinom{N}{k}\, p^k (1 - p)^{N - k}$ & $N p$ & $N p\,(1 - p)$ \\
  \rowcolor{fsfondo}
  \nomedist{Uniforme}{uniforme}{$M$ valori equiprobabili} & $0, 1, \ldots, M - 1$ & $\dfrac{1}{M}$ & $\dfrac{M -
    1}{2}$ & $\dfrac{M^2 - 1}{12}$ \\
  \nomedist{Poisson $\mathcal{P}(\lambda)$}{poisson}{eventi rari} & $0, 1, 2, \ldots$ & $\dfrac{\lambda^k}{k!}\,
    e^{-\lambda}$ & $\lambda$ & $\lambda$ \\
  \midrule
  \multicolumn{5}{l}{\color{fcolore}\Lato\bfseries Continue} \\
  \rowcolor{fsfondo}
  \nomedist{Uniforme $\mathcal{U}(a, b)$}{uniformec}{nessun punto privilegiato} & $(a, b)$ & $\dfrac{1}{b - a}$ &
    $\dfrac{a + b}{2}$ & $\dfrac{(b - a)^2}{12}$ \\
  \nomedist{Esponenziale $\mathcal{E}(\lambda)$}{esponenziale}{tempi di attesa} & $x \geq 0$ & $\lambda\,
    e^{-\lambda x}$ & $\dfrac{1}{\lambda}$ & $\dfrac{1}{\lambda^2}$ \\
  \rowcolor{fsfondo}
  \nomedist{Laplace $\mathcal{L}(\lambda)$}{laplacevar}{due esponenziali simmetrici} & $\R$ & $\dfrac{\lambda}{2}\,
    e^{-\lambda \abs{x}}$ & 0 & $\dfrac{2}{\lambda^2}$ \\
  \nomedist{Gaussiana $\mathcal{N}(\mu, \sigma^2)$}{gaussgen}{somme di tanti contributi} & $\R$ &
    $\dfrac{1}{\sqrt{2\pi\sigma^2}}\, e^{-\frac{(x - \mu)^2}{2\sigma^2}}$ & $\mu$ & $\sigma^2$ \\
  \rowcolor{fsfondo}
  \nomedist{Cauchy $\mathcal{C}(a, b)$}{cauchyvar}{code pesanti} & $\R$ & $\dfrac{1}{b\pi}\, \dfrac{1}{1 + \left(
    \frac{x - a}{b} \right)^2}$ & non esiste & non esiste \\
  \nomedist{Arcoseno}{arcoseno}{sinusoide con fase a caso} & $(-A, A)$ & $\dfrac{1}{\pi \sqrt{A^2 - x^2}}$ & 0 &
    $\dfrac{A^2}{2}$ \ \vv{sinusoideproc} \\
  \bottomrule
\end{tabular}}

\vspace{4mm}
Due regole permettono di adattare la tabella. Spostando e cambiando scala, $Y = aX + b$ ha media $a\,E[X] + b$ e
varianza $a^2\, \mathrm{VAR}(X)$ (\vv{lineare}): per esempio l'uniforme su $1, \ldots, 6$ (un dado) ha media
$\frac{M - 1}{2} + 1 = 3{,}5$ e la stessa varianza, $\frac{35}{12}$. E il valore quadratico medio è varianza più media
al quadrato, $E[X^2] = \sigma_X^2 + \mu_X^2$ (\vv{varrms}).
